\section{Introduction}

Periodic-orbit zeta functions for hyperbolic systems are built from additive
orbit data: a H\"older potential is summed along each primitive cycle and the
resulting weights enter a transfer operator or Euler product
\citep{Bowen1975,ParryPollicott1990}.  This additivity is not cosmetic.  It is
the interface that permits thermodynamic formalism to control an all-orbit
series.

The certified H6 H\'enon survivor studied in the preceding packages has a
mixing four-state coding and an intrinsic instability roof
\citep{WangP31}.  For each primitive cycle \(\gamma\), HCS-P54 split the Mahler
height of its algebraic return multiplier as
\begin{equation}\label{eq:split-intro}
 \mathcal H_\gamma=\ell_\gamma+\Egal_\gamma,
 \qquad \Egal_\gamma\geq 0.
\end{equation}
The physical length \(\ell_\gamma\) already has a source-backed pressure pole.
The full amplitude would inherit the weighted zeta machinery if one H\"older
observable had periodic sums \(\Egal_\gamma\), but that realization was left
open \citep{WangP54}.  The present paper subjects its most local forms to an
exact adversarial test.

Our first contribution is a cycle-incidence compiler.  A potential depending
on \(r\) consecutive states pairs linearly with the cyclic \(r\)-block count
of every orbit.  Any integer relation among those count vectors therefore
forces the same relation among the proposed orbit totals.  On the frozen H6
graph, the first three-block relation uses periods three, four, four and five.

The second contribution is new exact algebra for the required period-five
orbit.  Elimination produces a degree-six polynomial for its monodromy trace
and a reciprocal irreducible degree-twelve polynomial for its unstable
multiplier.  All six trace conjugates are real and lie outside \([-2,2]\), so
the Galois excess is an explicit sum of five positive reciprocal-pair lengths.
A rational Sturm interval already makes this excess too large for the
three-block relation.

The resulting theorem is deliberately finite-memory:
\begin{quote}
No locally constant potential depending on at most three consecutive H6
states realizes \(\Egal_\gamma\) on every primitive cycle.
\end{quote}
It does not refute an arbitrary H\"older potential.  Indeed, a finite family
of disjoint periodic orbits can always be interpolated by a sufficiently
long cylinder function.  We prove that finite-interpolation firewall and then
derive the correct one-sided all-period target: every forward high-block
incidence relation must have exponentially small Galois-excess discrepancy
if a H\"older realization on the one-sided H6 presentation exists.  The
unrestricted two-sided question remains open unless one supplies a
future-dependent cohomological reduction.  This quantitative condition is
the reusable bridge to the next paper.

No rational primes or Riemann zeros enter the construction.  The result
strengthens a local obstruction inside Route A; it does not construct a full
Galois-weighted determinant, a trace formula, or a Hilbert--P\'olya operator.
